\section{Infrastructure Requirements}
\label{sec:infrastructure}

% GUIDANCE: Cover the ground and launch infrastructure needed to build and
% operate the constellation.

% 6.1 Earth Ground Segment
%   - Optical ground stations with adaptive optics for atmospheric compensation.
%   - 4--8 geographically distributed sites for >99.5% availability.
%   - High-altitude, dry climates preferred (Atacama, Mauna Kea, La Palma).
%   - Multiple tracking telescopes per site for simultaneous ring links.

% 6.2 Launcher Production and Fleet
%   - 2,200--3,800 sats at 20/launch = 110--190 Starship launches.
%   - At weekly cadence: 2--4 years of dedicated launches.
%   - Realistic timeline sharing manifest: 5--8 years for full constellation.
%   - Fleet: 2--3 Starship vehicles with ~2-week turnaround.

% 6.3 Fuel Production
%   - ~3,400 tonnes propellant per Starship launch.
%   - 190 launches = ~650,000 tonnes total.
%   - Scaled Starbase facility can handle this.
%   - Lunar propellant depot option: water ice from PSRs for cislunar refueling.

% 6.4 Mars Ground Segment
%   - Mars thin atmosphere (1%) = negligible optical scintillation.
%   - Single optical terminal at settlement initially, with backup.
%   - Partial Mars co-orbit ring (~30 degrees) bridges outermost relay ring
%     to ground terminals.
%   - Additional terminals as settlement grows.
%
% Target length: ~1.5 pages.
